\section{The Vigil of the Nativity: Today You Shall Know}\label{sec:vigil}

On the morning of 24 December, Advent reaches its last hours. The Vigil of the
Nativity of the Lord is a first-class vigil, outranking every feast and even
the Fourth Sunday of Advent when 24 December falls on a Sunday (RG n.~30~a: the
Vigil takes the Sunday's place with no commemoration of the Sunday). The Roman
station returns to St.~Mary Major, where the season began four weeks earlier.
Purple vestments stay on the altar through the daytime Mass, the \textit{Gloria
in excelsis} remains silent, and Advent closes strictly at None in the
afternoon (RG n.~71); at First Vespers of Christmas, everything turns to white
and gold.

Yet the morning Mass already trembles with what evening will bring. The Introit
takes the words Moses spoke to hungry Israel on the eve of the manna in the
wilderness (Ex 16:6--7) and turns them directly toward tonight's Child:
\textit{Hodie scietis, quia veniet Dominus, et salvabit nos: et mane videbitis
gloriam ejus}---``This day you shall know that the Lord will come, and save us:
and in the morning you shall see his glory.'' In Exodus, Israel waited in the
desert for bread from heaven; on Christmas Eve, the Church waits for the living
Bread to lie in Bethlehem. The Communion repeats Isaiah 40:5: ``And the glory
of the Lord shall be revealed, and all flesh together shall see, that the mouth
of the Lord hath spoken.''

In the Epistle (Romans 1:1--6), St.~Paul introduces the Gospel promised before
by the prophets concerning God's Son, made of the seed of David according to
the flesh. And the Vigil Collect---the prayer that holds Advent's two horizons
in a single breath---begs that we who welcome the Redeemer with gladness
tonight may face him without terror when he returns as Judge:

\begin{massbox}{Collect of the Vigil of the Nativity}
\textbf{Latin.} Deus, qui nos redemptionis nostrae annua exspectatione
laetificas: praesta; ut Unigenitum tuum, quem Redemptorem laeti suscipimus,
venientem quoque judicem securi videamus, Dominum nostrum Jesum Christum,
Filium tuum: Qui tecum vivit et regnat in unitate.

\textbf{English.} O God, Who dost gladden us year by year with the expectation
of our redemption, grant that we, who now with joy receive Thine only begotten
Son as our Redeemer, may behold Him also without fear, when He cometh as our
judge, our Lord Jesus Christ. Who with thee.
\end{massbox}

\subsection{St.~Joseph's Silence and the Name of Jesus}

The Vigil Gospel is Matthew 1:18--21, the hidden crisis before the road to
Bethlehem. Mary, espoused to Joseph, is found with child of the Holy Ghost
before they came together: ``Whereupon Joseph her husband, being a just man,
and not willing publicly to expose her, was minded to put her away privately''
(Matt 1:19).

At Vigil Matins, all three lessons belong to St.~Jerome's commentary on this
Gospel (RGBR n.~222~a). St.~Jerome pauses first over why the Son of God was
conceived of an espoused virgin rather than an unwed one: ``First, that by the
descent of Joseph, Mary's family might be made known; secondly, that she might
not be stoned by the Jews as an adulteress; thirdly, that in her flight into
Egypt she might have the comfort of a husband,'' adding St.~Ignatius of
Antioch's fourth reason, ``that His birth might be hid from the Devil.''

Then Jerome turns to the harder verse: how can Scripture call Joseph ``just''
when he decides to send Mary away quietly rather than denounce her under the
Law? In his \textit{Commentary on Matthew} (Book~I, preserved in the
\textit{Catena Aurea}), Jerome gives an answer that opens the whole interior
stillness of St.~Joseph:

\begin{quote}
Or this may be considered a testimony to Mary, that Joseph, confident in her
purity, and wondering at what had happened, covered in silence that mystery
which he could not explain.
\end{quote}

Joseph does not suspect Mary of sin; he knows her holiness and stands in awe
before an act of God too large for him to grasp, unwilling to intrude where
heaven has acted. Into that reverent silence enters the angel in a dream:
``Joseph, son of David, fear not to take unto thee Mary thy wife, for that
which is conceived in her, is of the Holy Ghost. And she shall bring forth a
son: and thou shalt call his name Jesus. For he shall save his people from
their sins'' (Matt 1:20--21). St.~Jerome points to the Hebrew root of that
name: ``Jesus is a Hebrew word, meaning Saviour.'' And St.~Peter Chrysologus,
also gathered by St.~Thomas in the \textit{Catena Aurea} on Matthew 1:18,
marvels at the stainless order of that marriage: ``where an Angel is mediator,
faith is bridesmaid, where chastity is the giving away, virtue the gift,
conscience the judge, God the cause; where the conception is inviolateness, the
birth virginity, and the mother a virgin.''

\subsection{The Proclamation of the Ages}

In choir at Prime on 24 December, the Roman Martyrology announces the morrow
with a solemn chant known as the \textit{Kalenda}---from its opening words,
\textit{Octavo Kalendas Januarii}, the eighth day before the Kalends of
January. Recited in the Church's older ecclesiastical chronology
(Appendix~\ref{app:about}), the reader counts down the whole long patience of
the world: five thousand one hundred and ninety-nine years from the creation of
heaven and earth, two thousand nine hundred and fifty-seven from the flood, two
thousand and fifteen from the birth of Abraham, one thousand five hundred and
ten from Moses and the Exodus, one thousand and thirty-two from the anointing
of David, the sixty-fifth week of Daniel, the hundred and ninety-fourth
Olympiad, seven hundred and fifty-two years from the founding of Rome, the
forty-second year of the empire of Octavian Augustus, while the whole world was
at peace, in the sixth age of the world. Then the reader raises his voice and
all kneel:

\begin{quote}
Jesus Christus aeternus Deus, aeternique Patris Filius, mundum volens adventu
suo piissimo consecrare, de Spiritu sancto conceptus, novemque post
conceptionem decursis mensibus, \ldots{} in Bethlehem Judae nascitur ex Maria
Virgine factus Homo.
\end{quote}

``Jesus Christ, eternal God and Son of the eternal Father, willing to
consecrate the world by his most gracious coming, conceived of the Holy Ghost,
and nine months having passed since his conception, in Bethlehem of Juda is
born of the Virgin Mary, being made man.''

Preaching to his monks on that very sentence of the Martyrology in his
\textit{First Sermon for the Vigil of the Nativity}, St.~Bernard of Clairvaux
could scarcely move past its words:

\begin{quote}
Jesus Christ, the Son of God, is born in Bethlehem of Juda. What heart so stony
as not to be softened at these words? What soul is not melted at this voice of
her Beloved? What announcement could be sweeter? what intelligence more
enrapturing? Was its like ever heard before? or when did the world ever receive
such tidings? Jesus Christ, the Son of God, is born in Bethlehem of Juda. O
short word, telling of the Eternal Word abbreviated for us!
\end{quote}

As afternoon fades toward First Vespers on 24 December, the work of Advent is
done. Clear the table, finish the errands early if you can, leave room for
silence before the evening bells, and let the four verses of the French Advent
prose \textit{Rorate caeli} (Appendix~\ref{app:prayers}) have their final
hearing before nightfall: ``Be comforted, be comforted, my people; thy
salvation shall speedily come: why wilt thou waste away in sadness? why hath
sorrow seized thee? I will save thee; fear not: for I am the Lord thy God, the
Holy One of Israel, thy Redeemer.''
